\chapter{Нетривиальное KMS-состояние и селекция относительных скоростей}
\label{ch:v8-kms-nontracial-relative-rate-selector-gate}

\section{Проверяемый вопрос}

Полный генератор предыдущей главы примитивен, но содержит шесть
положительных относительных весов. Естественный кандидат селектора ---
подробный баланс относительно верного нетривиального состояния
\begin{equation}
 \rho_{a,b}=aI_{11}\oplus bI_{10},
 \qquad a,b>0,
 \qquad 11a+10b=1.
 \label{eq:v8-central-kms-state}
\end{equation}
Если такое состояние существует для полного процесса, отношение $b/a$
могло бы задать предпочтительное направление межуглового обмена.

\section{Точный запрет внутри текущего генератора}

KMS подробный баланс требует стационарности опорного состояния. Но уже
собранная полугруппа примитивна и имеет единственное верное стационарное
состояние
\begin{equation}
 \rho_*=\frac{I_{21}}{21}.
 \label{eq:v8-unique-tracial-kms-state}
\end{equation}
Следовательно,
\begin{equation}
 \rho_{a,b}\ \hbox{стационарно}
 \quad\Longleftrightarrow\quad
 a=b=\frac1{21}.
 \label{eq:v8-central-kms-only-trace}
\end{equation}

Машинная проверка независимо воспроизводит этот вывод. На концевой алгебре
KMS- и GNS-метрики центрального состояния представлены матрицей
\begin{equation}
 G_{a,b}=aI_{121}\oplus bI_{100}.
\end{equation}
В скане $10^{-4}\leq b/a\leq10^4$ одновременно стационарность и условие
$G\mathcal L=\mathcal L^*G$ проходят только при $b/a=1$.

Подбор положительных весов не помогает. Связывающие, $QLYR$- и $XLdR$-
переносы действуют на единицу исходного угла положительными матрицами с
суммарными следами
\begin{equation}
 13,\qquad 6,\qquad 6.
 \label{eq:v8-positive-transfer-traces}
\end{equation}
Их положительная линейная комбинация не может обнулиться, поэтому
взвешенная симметрия снова требует $a=b$.

\section{Что потребовалось бы для нетривиального KMS-отношения}

Пусть модульный гамильтониан имеет центральный перепад
\begin{equation}
 H_\Delta=0\cdot P_{11}+\Delta P_{10}.
\end{equation}
Нынешний самосопряжённый скачок имеет вид $J=V+V^*$, где $V$ направлен из
$\mathbb C^{11}$ в $\mathbb C^{10}$. При $\Delta\neq0$ он не является одной
боровской модой:
\begin{equation}
 [H_\Delta,V]=\Delta V,
 \qquad
 [H_\Delta,V^*]=-\Delta V^*.
 \label{eq:v8-directed-bohr-split}
\end{equation}
Нетривиальная KMS-динамика поэтому требует заменить каждый двусторонний
скачок двумя направленными каналами. Тогда подробный баланс даёт лишь
условное отношение
\begin{equation}
 \frac{\gamma_\uparrow}{\gamma_\downarrow}
 =e^{-\beta\Delta}
 =\frac ba.
 \label{eq:v8-kms-directed-rate-ratio}
\end{equation}
Но ни $\beta\Delta$, ни само состояние концевого сектора $\rho_{a,b}$ текущим
родителем не выбраны. Формула переносит свободу из отношения скоростей в
модульный разрыв. Базовые интенсивности трёх семейств переноса, трёх
калибровочных семейств и общий масштаб времени также остаются независимыми.

Ранние аффинные Гиббс-состояния Тома IV живут в семейных $M_3$-углах. В
проекте нет канонического отображения, поднимающего их в состояние
$M_{11}\oplus M_{10}$, поэтому использовать их здесь как готовую тепловую
плотность нельзя.

\begin{theorem}[Единственность следового KMS-состояния текущего процесса]
Для полного унитального примитивного генератора Тома VIII всякое верное
состояние, относительно которого выполняется KMS подробный баланс, обязано
совпадать с $I_{21}/21$. Следовательно, нетривиальное KMS-состояние не
выбирает относительные скорости шести уже построенных слагаемых.
\end{theorem}

\begin{nogocriterion}[Запрет скрытой подстановки модульного разрыва]
Нельзя объявлять $e^{-\beta\Delta}$ выведенным отношением скоростей, пока
сам оператор $H_\Delta$, его масштаб и направленное разбиение скачков не
получены из родителя. Выбор $\beta\Delta$ по желаемой скорости является
перепараметризацией подгонки.
\end{nogocriterion}

\begin{protocol}
\begin{enumerate}
 \item примитивность исходного процесса: \textbf{унаследована};
 \item единственное стационарное состояние: \textbf{$I_{21}/21$};
 \item нетривиальное центральное KMS-состояние: \textbf{невозможно};
 \item скан отношения $b/a$: \textbf{$10^{-4}\ldots10^4$};
 \item единственная проходящая точка: \textbf{$b/a=1$};
 \item обход подбором положительных весов: \textbf{невозможен};
 \item самосопряжённый скачок при $\Delta\neq0$ является одной боровской
 модой: \textbf{нет};
 \item направленное разбиение $V,V^*$: \textbf{математически возможно};
 \item отношение $\gamma_\uparrow/\gamma_\downarrow$: \textbf{условно
 равно $e^{-\beta\Delta}$};
 \item модульный разрыв $\beta\Delta$: \textbf{не выведен};
 \item селекция шести относительных скоростей: \textbf{не получена};
 \item следующий тест: \textbf{происхождение модульного боровского
 гамильтониана из родителя}.
\end{enumerate}
\end{protocol}

Машинный сертификат сохранён в
\path{s2t_v8_kms_nontracial_relative_rate_selector_gate_results.json}.

\begin{thebibliography}{9}
\bibitem{v8-kms-fagnola-umanita}
F.~Fagnola, V.~Umanit\`a,
\emph{Generators of detailed balance quantum Markov semigroups},
Infinite Dimensional Analysis, Quantum Probability and Related Topics
\textbf{10} (2007), 335--363; arXiv:0707.2147.
\bibitem{v8-kms-carlen-maas}
E.~A.~Carlen, J.~Maas,
\emph{Gradient flow and entropy inequalities for quantum Markov semigroups
with detailed balance}, Journal of Functional Analysis \textbf{273}
(2017), 1810--1869; arXiv:1609.01254.
\bibitem{v8-kms-temme}
K.~Temme, M.~J.~Kastoryano, M.~B.~Ruskai, M.~M.~Wolf, F.~Verstraete,
\emph{The $\chi^2$-divergence and mixing times of quantum Markov
processes}, Journal of Mathematical Physics \textbf{51} (2010), 122201;
arXiv:1005.2358.
\end{thebibliography}